\documentclass[12pt]{article}
\usepackage[UTF8]{ctex}
\usepackage{amsmath,amssymb}
\usepackage{geometry}
\usepackage{graphicx}
\usepackage{hyperref}
\usepackage{booktabs}
\geometry{margin=1in}

\newcommand{\atantwo}{\operatorname{atan2}}

\begin{document}

\title{EE5346 -- 自主机器人导航（2026年春季）\\ 作业\#5\\ MLE 与扩展卡尔曼滤波器}
\author{}
\date{}
\maketitle

%==============================================================
\section*{问题 1：最大似然估计 (MLE)}

测量方程为
\[
z_k = x + n_k,\qquad n_k \sim \mathcal{N}(0,\sigma^2),\qquad k=1,2,\dots,N,
\]
其中 $x$ 为待估计的机器人位置（未知常数），$n_k$ 为零均值高斯噪声（题目中的 $\mathcal{N}(0,\sigma)$ 中 $\sigma$ 视为标准差，方差为 $\sigma^2$），$N$ 个测量相互独立。由 $z_k\mid x \sim \mathcal{N}(x,\sigma^2)$，单测量的概率密度函数为
\[
p(z_k\mid x)=\frac{1}{\sqrt{2\pi\sigma^2}}\exp\!\left[-\frac{(z_k-x)^2}{2\sigma^2}\right].
\]

\subsection*{似然函数 $L(x)=p(Z\mid x)$}

由各测量独立，联合 PDF 等于各 PDF 的乘积：
\[
L(x)=\prod_{k=1}^{N}p(z_k\mid x)
=\left(\frac{1}{\sqrt{2\pi\sigma^2}}\right)^{\!N}
\exp\!\left[-\frac{1}{2\sigma^2}\sum_{k=1}^{N}(z_k-x)^2\right].
\]
即
\[
\boxed{\;
L(x)=(2\pi\sigma^2)^{-N/2}\exp\!\left[-\frac{1}{2\sigma^2}\sum_{k=1}^{N}(z_k-x)^2\right].
\;}
\]

\subsection*{负对数似然 $\mathcal{L}(x)=-\ln L(x)$}

利用 $\ln(ab)=\ln a+\ln b$ 与 $\ln(\mathrm{e}^c)=c$：
\[
\boxed{\;
\mathcal{L}(x)=\frac{N}{2}\ln(2\pi\sigma^2)+\frac{1}{2\sigma^2}\sum_{k=1}^{N}(z_k-x)^2.
\;}
\]
第一项与 $x$ 无关，因此
\[
x^{*}=\arg\min_{x}\mathcal{L}(x)\;\equiv\;\arg\min_{x}\sum_{k=1}^{N}(z_k-x)^2,
\]
即在高斯噪声下，MLE 等价于最小二乘问题。

\subsection*{求解 MLE $x^{*}$}

对 $\mathcal{L}(x)$ 求导并令其为零：
\[
\frac{d\mathcal{L}}{dx}=\frac{1}{\sigma^2}\sum_{k=1}^{N}(x-z_k)=0
\;\Longrightarrow\;
Nx-\sum_{k=1}^{N}z_k=0.
\]
二阶导数 $\dfrac{d^{2}\mathcal{L}}{dx^{2}}=\dfrac{N}{\sigma^{2}}>0$，确认为极小值。因此
\[
\boxed{\;x^{*}=\frac{1}{N}\sum_{k=1}^{N}z_k.\;}
\]
即\textbf{在 i.i.d.\ 高斯噪声下，机器人位置的 MLE 等于所有测量值的样本均值}。

\subsection*{数值结果}

代入 $N=3$，$z_1=3,\;z_2=4,\;z_3=7$：
\[
\boxed{\;x^{*}=\frac{3+4+7}{3}=\frac{14}{3}\approx 4.667.\;}
\]

\newpage
%==============================================================
\section*{问题 2：教材式 (8.32) 的推导}

考虑非线性运动模型（加性过程噪声形式）
\[
\mathbf{x}_k=f(\mathbf{x}_{k-1},\mathbf{u}_k)+\mathbf{w}_k,
\qquad \mathbf{w}_k\sim\mathcal{N}(\mathbf{0},\mathbf{R}_k),
\]
其中 $\mathbf{w}_k$ 与历史状态独立。上一时刻的后验为
\[
\mathbf{x}_{k-1}\sim\mathcal{N}(\hat{\mathbf{x}}_{k-1},\hat{\mathbf{P}}_{k-1}).
\]

\subsection*{在估计点处一阶泰勒展开}

由于 $f$ 非线性，EKF 在 $\hat{\mathbf{x}}_{k-1}$ 处线性化：
\[
f(\mathbf{x}_{k-1},\mathbf{u}_k)\approx
f(\hat{\mathbf{x}}_{k-1},\mathbf{u}_k)
+\mathbf{F}\,(\mathbf{x}_{k-1}-\hat{\mathbf{x}}_{k-1}),
\qquad
\mathbf{F}=\left.\frac{\partial f}{\partial \mathbf{x}}\right|_{\hat{\mathbf{x}}_{k-1},\mathbf{u}_k}.
\]
代入状态方程：
\begin{equation}
\mathbf{x}_k\approx
f(\hat{\mathbf{x}}_{k-1},\mathbf{u}_k)
+\mathbf{F}\,(\mathbf{x}_{k-1}-\hat{\mathbf{x}}_{k-1})
+\mathbf{w}_k.
\tag{$\star$}
\end{equation}

\subsection*{预测均值 $\tilde{\mathbf{x}}_k$}

由\textbf{均值（向量函数期望）的定义} $\tilde{\mathbf{x}}_k=E[\mathbf{x}_k]$，对 $(\star)$ 两边取期望并利用线性性：
\[
\tilde{\mathbf{x}}_k=
f(\hat{\mathbf{x}}_{k-1},\mathbf{u}_k)
+\mathbf{F}\underbrace{E[\mathbf{x}_{k-1}-\hat{\mathbf{x}}_{k-1}]}_{=\mathbf{0}}
+\underbrace{E[\mathbf{w}_k]}_{=\mathbf{0}}.
\]
故
\[
\boxed{\;\tilde{\mathbf{x}}_k=f(\hat{\mathbf{x}}_{k-1},\mathbf{u}_k).\;}
\]

\subsection*{预测协方差 $\tilde{\mathbf{P}}_k$}

记误差 $\delta\mathbf{x}_{k-1}=\mathbf{x}_{k-1}-\hat{\mathbf{x}}_{k-1}$。由 $(\star)$ 及上一节结果：
\[
\mathbf{x}_k-\tilde{\mathbf{x}}_k=\mathbf{F}\,\delta\mathbf{x}_{k-1}+\mathbf{w}_k.
\]
按\textbf{协方差矩阵的定义}
$\tilde{\mathbf{P}}_k=E\!\left[(\mathbf{x}_k-\tilde{\mathbf{x}}_k)(\mathbf{x}_k-\tilde{\mathbf{x}}_k)^\mathrm{T}\right]$
展开：
\[
\tilde{\mathbf{P}}_k=
E\!\left[
\mathbf{F}\delta\mathbf{x}_{k-1}\delta\mathbf{x}_{k-1}^\mathrm{T}\mathbf{F}^\mathrm{T}
+\mathbf{F}\delta\mathbf{x}_{k-1}\mathbf{w}_k^\mathrm{T}
+\mathbf{w}_k\delta\mathbf{x}_{k-1}^\mathrm{T}\mathbf{F}^\mathrm{T}
+\mathbf{w}_k\mathbf{w}_k^\mathrm{T}
\right].
\]
由于 $\mathbf{w}_k$ 与 $\mathbf{x}_{k-1}$ 独立，
$E[\delta\mathbf{x}_{k-1}\mathbf{w}_k^\mathrm{T}]=\mathbf{0}$，交叉项消失；又
\[
E[\delta\mathbf{x}_{k-1}\delta\mathbf{x}_{k-1}^\mathrm{T}]=\hat{\mathbf{P}}_{k-1},
\qquad
E[\mathbf{w}_k\mathbf{w}_k^\mathrm{T}]=\mathbf{R}_k,
\]
因此
\[
\boxed{\;
\tilde{\mathbf{P}}_k=\mathbf{F}\,\hat{\mathbf{P}}_{k-1}\,\mathbf{F}^\mathrm{T}+\mathbf{R}_k.
\;}
\]

\subsection*{关于非加性噪声的说明}

若噪声以 $f(\mathbf{x}_{k-1},\mathbf{u}_k,\mathbf{w}_k)$ 方式进入模型（如本作业第 3 题中过程噪声乘以输入耦合矩阵），则需引入噪声雅可比
\[
\mathbf{G}=\left.\frac{\partial f}{\partial \mathbf{w}}\right|_{\hat{\mathbf{x}}_{k-1},\mathbf{u}_k,\mathbf{0}},
\]
此时预测协方差为
\[
\tilde{\mathbf{P}}_k=\mathbf{F}\,\hat{\mathbf{P}}_{k-1}\,\mathbf{F}^\mathrm{T}
+\mathbf{G}\,\mathbf{Q}\,\mathbf{G}^\mathrm{T}.
\]
教材式 (8.32) 通过将 $\mathbf{G}\mathbf{Q}\mathbf{G}^\mathrm{T}$ 整体记为 $\mathbf{R}_k$，与上述加性形式一致。

\newpage
%==============================================================
\section*{问题 3：EKF 雅可比矩阵 $\mathbf{F}$ 与 $\mathbf{H}$}

记状态、输入、地标分别为
\[
\mathbf{x}_k=\begin{bmatrix}x_k\\ y_k\\ \theta_k\end{bmatrix},\quad
\mathbf{u}_k=\begin{bmatrix}v_k\\ \omega_k\end{bmatrix},\quad
\mathbf{y}^l=\begin{bmatrix}x_l\\ y_l\end{bmatrix}.
\]

\subsection*{运动模型 $\Rightarrow$ 雅可比 $\mathbf{F}$}

题目给定的运动模型为
\[
\mathbf{x}_k = \mathbf{x}_{k-1}
+ \Delta t
\begin{bmatrix}
\cos\theta_{k-1} & 0\\
\sin\theta_{k-1} & 0\\
0 & 1
\end{bmatrix}
(\mathbf{u}_k+\mathbf{w}_k),
\qquad \mathbf{w}_k\sim\mathcal{N}(\mathbf{0},\mathbf{Q}).
\]
取 $\mathbf{w}_k=\mathbf{0}$ 用于求 $\mathbf{F}$，逐分量展开：
\[
\begin{aligned}
x_k       &= x_{k-1} + \Delta t\,v_k\cos\theta_{k-1},\\
y_k       &= y_{k-1} + \Delta t\,v_k\sin\theta_{k-1},\\
\theta_k  &= \theta_{k-1} + \Delta t\,\omega_k.
\end{aligned}
\]

\subsubsection*{一阶近似形式}

在 $(\hat{\mathbf{x}}_{k-1},\mathbf{u}_k)$ 处展开：
\[
f(\mathbf{x}_{k-1},\mathbf{u}_k)\approx
f(\hat{\mathbf{x}}_{k-1},\mathbf{u}_k)
+\mathbf{F}\,(\mathbf{x}_{k-1}-\hat{\mathbf{x}}_{k-1}).
\]

\subsubsection*{求 $\mathbf{F}=\partial f/\partial \mathbf{x}_{k-1}$}

逐项求偏导：
\begin{align*}
\frac{\partial x_k}{\partial x_{k-1}} &= 1, &
\frac{\partial x_k}{\partial y_{k-1}} &= 0, &
\frac{\partial x_k}{\partial \theta_{k-1}} &= -\Delta t\,v_k\sin\theta_{k-1},\\[4pt]
\frac{\partial y_k}{\partial x_{k-1}} &= 0, &
\frac{\partial y_k}{\partial y_{k-1}} &= 1, &
\frac{\partial y_k}{\partial \theta_{k-1}} &= \Delta t\,v_k\cos\theta_{k-1},\\[4pt]
\frac{\partial \theta_k}{\partial x_{k-1}} &= 0, &
\frac{\partial \theta_k}{\partial y_{k-1}} &= 0, &
\frac{\partial \theta_k}{\partial \theta_{k-1}} &= 1.
\end{align*}

因此
\[
\boxed{\;
\mathbf{F}=
\begin{bmatrix}
1 & 0 & -\Delta t\,v_k\sin\hat{\theta}_{k-1}\\[4pt]
0 & 1 & \;\;\;\Delta t\,v_k\cos\hat{\theta}_{k-1}\\[4pt]
0 & 0 & 1
\end{bmatrix}.
\;}
\]
矩阵在估计点 $\hat{\theta}_{k-1}$ 与里程计输入 $v_k$ 处求值。

\subsection*{测量模型 $\Rightarrow$ 雅可比 $\mathbf{H}$}

\subsubsection*{中间变量}

定义
\[
\Delta x = x_l-x_k-d\cos\theta_k,\qquad
\Delta y = y_l-y_k-d\sin\theta_k,
\]
\[
r=\sqrt{\Delta x^2+\Delta y^2},\qquad r^2=\Delta x^2+\Delta y^2.
\]
测量模型为
\[
h(\mathbf{x}_k,\mathbf{y}^l)=
\begin{bmatrix}
r\\[3pt]
\atantwo(\Delta y,\Delta x)-\theta_k
\end{bmatrix}.
\]

\subsubsection*{一阶近似形式}

在预测位姿 $\tilde{\mathbf{x}}_k$ 处展开：
\[
h(\mathbf{x}_k,\mathbf{y}^l)\approx
h(\tilde{\mathbf{x}}_k,\mathbf{y}^l)
+\mathbf{H}\,(\mathbf{x}_k-\tilde{\mathbf{x}}_k).
\]

\subsubsection*{辅助偏导}

\begin{align*}
\frac{\partial \Delta x}{\partial x_k}&=-1, &
\frac{\partial \Delta y}{\partial y_k}&=-1,\\[4pt]
\frac{\partial \Delta x}{\partial \theta_k}&=d\sin\theta_k, &
\frac{\partial \Delta y}{\partial \theta_k}&=-d\cos\theta_k.
\end{align*}

\subsubsection*{(a) 距离 $r$ 部分}

\[
\frac{\partial r}{\partial x_k}=\frac{1}{r}\!\left(\Delta x\cdot(-1)+\Delta y\cdot 0\right)=-\frac{\Delta x}{r},
\qquad
\frac{\partial r}{\partial y_k}=-\frac{\Delta y}{r},
\]
\[
\frac{\partial r}{\partial \theta_k}
=\frac{1}{r}\!\left(\Delta x\cdot d\sin\theta_k+\Delta y\cdot(-d\cos\theta_k)\right)
=\frac{d}{r}\,(\Delta x\sin\theta_k-\Delta y\cos\theta_k).
\]

\subsubsection*{(b) 方位 $\phi=\atantwo(\Delta y,\Delta x)-\theta_k$ 部分}

利用
\[
\frac{\partial\,\atantwo(v,u)}{\partial u}=-\frac{v}{u^2+v^2},\qquad
\frac{\partial\,\atantwo(v,u)}{\partial v}=\frac{u}{u^2+v^2},
\]
通过链式法则得：
\begin{align*}
\frac{\partial\phi}{\partial x_k}
&=-\frac{\Delta y}{r^2}\cdot(-1)+\frac{\Delta x}{r^2}\cdot 0
=\frac{\Delta y}{r^2},\\[6pt]
\frac{\partial\phi}{\partial y_k}
&=-\frac{\Delta y}{r^2}\cdot 0+\frac{\Delta x}{r^2}\cdot(-1)
=-\frac{\Delta x}{r^2},\\[6pt]
\frac{\partial\phi}{\partial\theta_k}
&=-\frac{\Delta y}{r^2}\cdot d\sin\theta_k
+\frac{\Delta x}{r^2}\cdot(-d\cos\theta_k)-1\\[4pt]
&=-\frac{d}{r^2}\!\left(\Delta x\cos\theta_k+\Delta y\sin\theta_k\right)-1.
\end{align*}

\subsubsection*{最终 $\mathbf{H}$（$2\times 3$ 矩阵）}

\[
\boxed{\;
\mathbf{H}=
\begin{bmatrix}
-\dfrac{\Delta x}{r} &
-\dfrac{\Delta y}{r} &
\;\;\dfrac{d}{r}\!\left(\Delta x\sin\theta_k-\Delta y\cos\theta_k\right)\\[14pt]
\;\;\dfrac{\Delta y}{r^2} &
-\dfrac{\Delta x}{r^2} &
-\dfrac{d}{r^2}\!\left(\Delta x\cos\theta_k+\Delta y\sin\theta_k\right)-1
\end{bmatrix},
\;}
\]
其中
\[
\Delta x=x_l-x_k-d\cos\theta_k,\quad
\Delta y=y_l-y_k-d\sin\theta_k,\quad
r=\sqrt{\Delta x^2+\Delta y^2}.
\]
所有项在预测位姿 $\tilde{\mathbf{x}}_k=(\tilde{x}_k,\tilde{y}_k,\tilde{\theta}_k)$ 与已知地标 $\mathbf{y}^l$ 处求值。

\newpage
%==============================================================
\section*{答案速览}

\begin{table}[h]
\centering
\renewcommand{\arraystretch}{1.5}
\begin{tabular}{cl}
\toprule
\textbf{题号} & \textbf{关键结果}\\
\midrule
1 & $L(x)=(2\pi\sigma^2)^{-N/2}\exp\!\left[-\dfrac{1}{2\sigma^2}\sum_k(z_k-x)^2\right]$ \\
1 & $\mathcal{L}(x)=\dfrac{N}{2}\ln(2\pi\sigma^2)+\dfrac{1}{2\sigma^2}\sum_k(z_k-x)^2$ \\
1 & $x^{*}=\dfrac{1}{N}\sum_k z_k=\dfrac{14}{3}\approx 4.667$ \\
2 & $\tilde{\mathbf{x}}_k=f(\hat{\mathbf{x}}_{k-1},\mathbf{u}_k)$ \\
2 & $\tilde{\mathbf{P}}_k=\mathbf{F}\hat{\mathbf{P}}_{k-1}\mathbf{F}^\mathrm{T}+\mathbf{R}_k$ \\
3 & $\mathbf{F}=\begin{bmatrix}1&0&-\Delta t\,v_k\sin\hat{\theta}_{k-1}\\0&1&\Delta t\,v_k\cos\hat{\theta}_{k-1}\\0&0&1\end{bmatrix}$ \\
3 & $\mathbf{H}=\begin{bmatrix}-\Delta x/r & -\Delta y/r & (d/r)(\Delta x\sin\theta_k-\Delta y\cos\theta_k)\\ \Delta y/r^2 & -\Delta x/r^2 & -(d/r^2)(\Delta x\cos\theta_k+\Delta y\sin\theta_k)-1\end{bmatrix}$ \\
\bottomrule
\end{tabular}
\end{table}

\end{document}
